% Source: https://artofproblemsolving.com/community/c3h1659773_a_problem_on_summations
% Stable identity: aops-topic-1659773-q01
\textbf{(a)} Each of the $n$ terms of $B_n$ is at least its last term $\dfrac{1}{\sqrt{2n}}$, so
\[
B_n\ge\dfrac{n}{\sqrt{2n}}=\sqrt{\dfrac {n}{2}},
\]
which grows without bound. Since $\dfrac{1}{\sqrt{2k-1}}>\dfrac{1}{\sqrt{2k}}$ term by term, $A_n>B_n$ grows without bound as well.

\textbf{(b)} Compare the sums term by term, shifting by one place. For every $k\ge1$,
\[
\dfrac{1}{\sqrt{2k}}<\dfrac{1}{\sqrt{2k-1}}\qquad\text{and, for }k\ge2,\qquad \dfrac{1}{\sqrt{2k-1}}<\dfrac{1}{\sqrt{2k-2}} .
\]
Summing the first inequality gives $B_n<A_n$. Summing the second over $k=2,\dots,n$ and adding the first term $1$ gives
\[
A_n<1+\dfrac{1}{\sqrt2}+\dfrac{1}{\sqrt4}+\cdots+\dfrac{1}{\sqrt{2n-2}}=1+B_{n-1}<1+B_n .
\]
Dividing by $B_n>0$,
\[
1<\dfrac{A_n}{B_n}<1+\dfrac{1}{B_n} .
\]
By (a), $\dfrac{1}{B_n}\to0$, so by squeezing,
\[
\lim_{n\to\infty}\dfrac{A_n}{B_n}=\boxed{1}.
\]
(The finite sums in the post, $a=A_4$ and $b=B_5$, are ordinary numbers, $a\approx2.4025$ and $b\approx2.2851$, with no special closed form for $\dfrac {a}{b}$.)
